\documentclass[12pt]{article}

% Ustawienia strony
\usepackage{lipsum}
\usepackage[margin=1.5cm]{geometry}
\usepackage{graphicx}
\usepackage{enumitem}
\usepackage{siunitx}
\usepackage{amsmath}
\usepackage{cancel}
\usepackage{makecell}
\usepackage{mathtools}
\usepackage{floatrow}
\usepackage{caption}
\usepackage[T1]{fontenc}
\usepackage{fancyvrb}
\usepackage{subcaption}
\usepackage{blindtext}
\usepackage{hyperref}
\usepackage{amssymb}
\usepackage{minted}
\usepackage{xcolor} % to access the named color LightGray
\definecolor{LightGray}{gray}{0.9}

\setminted{
	frame=lines,
	framesep=1mm,
	baselinestretch=1,
	bgcolor=LightGray,
	fontsize=\footnotesize,
	linenos
}

\hypersetup{
	colorlinks=true,
	linkcolor=blue,
	filecolor=magenta,      
	urlcolor=cyan,
	pdfpagemode=FullScreen,
}

\urlstyle{same}

\fvset{tabsize=4}
\setlist{leftmargin=15pt,labelindent=15pt}
\setlist[enumerate]{noitemsep,wide=0pt, leftmargin=*}
\setlength{\parindent}{0pt}

% Ustawienia stopki
\usepackage{lastpage}
\usepackage{fancyhdr}
\pagestyle{fancy}
\fancyhf{} % sets both header and footer to nothing
\renewcommand{\headrulewidth}{0pt}
\cfoot{\thepage\ z \pageref*{LastPage}}

\begin{document}
\textbf{Robert Mesek}\text{,} % Autor
\text{15.04.2025}\text{ } % Data

\begin{center}
	\section*{Optymalizacja Kodu na Różne Architektury \\
	  Zadanie 1 - Normalizacja Tekstu}
	\mbox{}
\end{center}

\section*{Cel zadania}
Celem tego zadania jest implementacja funkcji, która przetwarza duże ciągi
tekstowe (np. zawartość dużego pliku) wykonując szereg operacji normalizacyjnych 
oraz optymalizacyjnych. \\ \\
Funkcja powinna:
\begin{itemize}
	\item Usuwać znaki spoza zakresu drukowalnych \\
	(np. wszystkie znaki o kodzie poniżej 32 lub powyżej 126).
	\item Zamieniać sekwencje białych znaków (spacje, tabulatory, nowe linie) \\
	na pojedynczą spację.
	\item Konwertować wszystkie litery do małych liter.
	\item Konwertować interpunkcję na przecinki.
	\item Eliminować duplikaty wyrazów występujących jeden po drugim \\
	(np. "hello hello world" → "hello world").
\end{itemize}
\subsection*{Przykład}
\begin{minted}{text}
Source:     "Hello hello world... World."
Normalized: "hello  world,,, ,"

Source:     "\t\n  First.first, ,, second   second; third\n\n"
Normalized: " first,, ,, second , third"

Source:     "     leading space and trailing space     "
Normalized: " leading space and trailing space "

Source:     "word1 word1, word2 word2"
Normalized: "word1 , word2 "

Source:     "Hello\x07 World\x7F!"
Normalized: "hello world,"
\end{minted}

\newpage

\section*{Dane techniczne}
\subsection*{System Operacyjny}
\begin{minted}{text}
$ hostnamectl
Virtualization: wsl
Operating System: Ubuntu 24.04.2 LTS
Kernel: Linux 5.15.167.4-microsoft-standard-WSL2
\end{minted}

\subsection*{Kompilator}
\begin{minted}{text}
$ g++ --version
g++ (Ubuntu 13.3.0-6ubuntu2~24.04) 13.3.0
\end{minted}

\subsection*{Procesor}
\begin{minted}{text}
$ lscpu
Architecture:             x86_64
  CPU op-mode(s):         32-bit, 64-bit
  Address sizes:          48 bits physical, 48 bits virtual
  Byte Order:             Little Endian
CPU(s):                   24
  On-line CPU(s) list:    0-23
Vendor ID:                AuthenticAMD
  Model name:             AMD Ryzen 9 7900X 12-Core Processor
    CPU family:           25
    Model:                97
    Thread(s) per core:   2
    Core(s) per socket:   12
    Socket(s):            1
    Stepping:             2
    BogoMIPS:             9399.78
    Flags:                fpu vme de pse tsc msr pae mce cx8 apic sep mtrr pge mca cmov pat pse36 
                          clflush mmx fxsr sse sse2 ht syscall nx mmxext fxsr_opt pdpe1gb rdtscp 
                          lm constant_tsc rep_good nopl tsc_reliable nonsto p_tsc cpuid extd_apicid 
                          pni pclmulqdq ssse3 fma cx16 sse4_1 sse4_2 movbe popcnt aes xsave avx f16c 
                          rdrand hypervisor lahf_lm cmp_legacy svm cr8_legacy abm sse4a misalignsse 
                          3dnowprefetch osvw topoext perfctr_core ssbd ibrs ibpb stibp vmmcall 
                          fsgsbase bmi1 avx2 smep bmi2 erms invpcid avx512f avx512dq rdseed adx smap 
                          avx512ifma clflushopt clwb avx512cd sha_ni avx512bw avx512vl xsaveo pt 
                          xsavec xgetbv1 xsaves avx512_bf16 clzero xsaveerptr arat npt nrip_save 
                          tsc_scale vmcb_clean flushbyasid decodeassists pausefilter pfthreshold 
                          v_vmsave_vmload avx512vbmi umip avx512_vbmi2 gf ni vaes vpclmulqdq 
                          avx512_vnni avx512_bitalg avx512_vpopcntdq rdpid fsrm
Caches (sum of all):      
  L1d:                    384 KiB (12 instances)
  L1i:                    384 KiB (12 instances)
  L2:                     12 MiB (12 instances)
  L3:                     32 MiB (1 instance)
\end{minted}

\newpage

\section*{Obliczanie czasu oraz profilowanie}
Czas został zmierzony dla zadanej liczby powtórzeń przy pomocy funkcji
\begin{minted}{cpp}
const auto start = std::chrono::steady_clock::now();
for (int i = 0; i < repeats; ++i) {
	normalized_content = normalize_text(content);
}
const auto finish = std::chrono::steady_clock::now();
const std::chrono::duration<double> elapsed_seconds = finish - start;
\end{minted}

Fragment danych wejściowych (w zadaniu użyłem wersji o rozmiarach \\
\verb|100KB|: \verb|1000| powtórzeń, \verb|100MB|: \verb|1| powtórzenie)

\begin{minted}{text}
	fpCEAO 
	fpCEAOGDFl XBFFP } CueBMKk CueBMKk vRatVsQl 	
	biNpsXGYw ITr BqHiYYKNdq 

<?>KQSVANXP <?>WnDAMnLjs hyoBoKZ xosRuQ vNfnGT : <?>: sGZHP uuDnkhNiBy 	 
VSyL IPegMxiL MxiL lbhYaRAt lbhYaRAt CUVWEIFBDQ 
- FgFVuKt <?>-
\end{minted}

Do profilowania użyłem narzędzia \verb|perf|.











\section*{Wersja bazowa}
Implementacja polega na jednokrotnym przejściu przez tekst, mając zmienną
na wynikowy tekst oraz zmienną na aktualnie ostatni i przetwarzany wyraz.
Iteracja odbywa się litera po literze, dodając ją do odpowiedniej zmiennej
w zależności od stanu. 

\begin{minted}{cpp}
/**
 * @brief Normalizes a text string according to specified rules.
 *
 * Rules:
 * 1. Remove non-printable ASCII characters (codes < 32 or > 126).
 * 2. Replace sequences of whitespace characters with a single space.
 * 3. Convert all letters to lowercase.
 * 4. Convert punctuation to commas.
 * 5. Eliminate consecutive duplicate words ("hi hi world"->"hi world").
 *
 * @param text The input string to normalize.
 * @return The normalized string.
 */
 std::string normalize_text(const std::string& text) {
  std::string result;
  std::string word, last_word;
  enum State { NORMAL, COMMA, SPACE, DUPLICATE_SPACE } state = NORMAL;

  for (char c : text) {
    if (c < 32 || c > 126) continue;  // rule 1

    if (c == ' ') {
      if (state == SPACE || state == DUPLICATE_SPACE) {  // rule 2
        state = DUPLICATE_SPACE;
      } else {
        state = SPACE;
      }
    } else if (isalpha(c)) {  // rule 3
      c = tolower(c);
      state = NORMAL;
    } else if (ispunct(c)) {  // rule 4
      c = ',';
      state = COMMA;
    } else {
      state = NORMAL;
    }

    if (state == COMMA || state == SPACE) {
      // end of a word
      if (!word.empty()) {  // rule 5
        if (word != last_word) result += word;
        last_word = word;
        word.clear();
      }
      result += c;  // only delimiter
    } else if (state == NORMAL) {
      // in a word
      word += c;  // only alphanumeric
    }
  }
  if (word != last_word) result += word;  // last word

  return result;
}
\end{minted}

Czas działania dla plików o rozmiarze \verb|100KB|
\begin{minted}{text}
normalize1 -O0 [Benchmark (1000 iterations): 1.97064 seconds]
normalize1 -O2 [Benchmark (1000 iterations): 0.64529 seconds]
\end{minted}

Czas działania dla plików o rozmiarze \verb|100MB|
\begin{minted}{text}
normalize1 -O0 [100MB][Benchmark (1 iterations): 1.98192 seconds]
normalize1 -O2 [100MB][Benchmark (1 iterations): 0.710283 seconds]
\end{minted}

Wynik profilowania
\begin{minted}{text}
54.56%  a.out    libc.so.6             [.] __mcount_internal
23.67%  a.out    libc.so.6             [.] _mcount
 4.25%  a.out    a.out                 [.] normalize_text(std::__cxx11::basic_string<char, ...
 4.13%  a.out    a.out                 [.] bool __gnu_cxx::operator!=<char const*, ...
 2.57%  a.out    libstdc++.so.6.0.33   [.] std::__cxx11::basic_string<char, ...
 2.12%  a.out    a.out                 [.] __gnu_cxx::__normal_iterator<char const*, ...
 2.09%  a.out    a.out                 [.] __gnu_cxx::__normal_iterator<char const*, ...
 1.51%  a.out    a.out                 [.] __gnu_cxx::__normal_iterator<char const*, ...
 1.11%  a.out    a.out                 [.] std::__cxx11::basic_string<char, ...
 1.10%  a.out    libc.so.6             [.] tolower
 0.61%  a.out    a.out                 [.] bool std::operator==<char, std::char_traits<char>, ...
 0.50%  a.out    libstdc++.so.6.0.33   [.] std::__cxx11::basic_string<char, ...
 0.29%  a.out    a.out                 [.] bool std::operator!=<char, std::char_traits<char>, ...
\end{minted}

Z analizy wyników możemy zaobserwować, że czas wykonania funkcji wygląda przyzwoicie,
jedną z rzeczy, którą możemy spróbować zredukować, która nie jest wymaga w algorytmie, 
jest dużo operacji na iteratorze. \\

Wynik ten jednak widzimy, że jest znacznie poprawiony przy ustawieniu optymalizacji
\verb|-O2|, co może sugerować, że kompilator mógł zoptymalizować operacje na iteratorze. \\

Kolejną obserwacją jest liniowa złożoność czasowa, co jest zgodne z oczekiwaniami, ponieważ
przechodzimy przez tekst tylko raz. \\

Co do zużycia pamięci, jest ona akceptowalna, ponieważ trzymamy tylko 2 zmienne pomocnicze
ograniczone do rozmiaru wyrazu oraz wynikowy tekst, który jest potrzebny, ponieważ pozwala
nam nie niszczyć oryginalnego tekstu oraz pozwala na efektywną obsługę powtarzających się 
wyrazów, co było by problematyczne przy usuwaniu ich w miejscu wymagającym przesunięcia
wszystkich pozostałych wyrazów w takiej strukturze danych. \\






\section*{Zwyczajna pętla oraz prealokacja}
W tej wersji zamieniono pętlę \verb|range-for| na pętlę \verb|for| oraz przetestowano
prealokację zmiennych \verb|result|, \verb|word| oraz \verb|last_word|. \\
\begin{minted}{cpp}
result.reserve(text.length());
word.reserve(256);
last_word.reserve(256);
\end{minted}
Użycie prealokacji zmiennych \textbf{nie przyniosło żadnej mierzalnej poprawy} w przypadku 
1000 iteracji po pliku o rozmiarze \verb|100KB| oraz 1 iteracji po pliku o rozmiarze
\verb|100MB|. \\

\begin{minted}{cpp}
std::string normalize_text(const std::string& text) {
  std::string result;
  std::string word, last_word;
  enum State { NORMAL, COMMA, SPACE, DUPLICATE_SPACE } state = NORMAL;

  for (size_t i = 0; i < text.length(); ++i) {
    char c = text[i];
    if (c < 32 || c > 126) continue;  // rule 1

    if (c == ' ') {
      if (state == SPACE || state == DUPLICATE_SPACE) {  // rule 2
        state = DUPLICATE_SPACE;
      } else {
        state = SPACE;
      }
    } else if (isalpha(c)) {  // rule 3
      c = tolower(c);
      state = NORMAL;
    } else if (ispunct(c)) {  // rule 4
      c = ',';
      state = COMMA;
    } else {
      state = NORMAL;
    }

    if (state == COMMA || state == SPACE) {
      // end of a word
      if (!word.empty()) {  // rule 5
        if (word != last_word) result += word;
        last_word = std::move(word);
        word.clear();
      }
      result += c;  // only delimiter
    } else if (state == NORMAL) {
      // in a word
      word += c;  // only alphanumeric
    }
  }
  if (word != last_word) result += word;  // last word

  return result;
}
\end{minted}

Czas działania dla plików o rozmiarze \verb|100KB|
\begin{minted}{text}
normalize2 -O0 [Benchmark (1000 iterations): 1.31727 seconds]
normalize2 -O2 [Benchmark (1000 iterations): 0.635332 seconds]
\end{minted}

Wynik profilowania
\begin{minted}{text}
34.90%  a.out    a.out                 [.] normalize_text(std::__cxx11::basic_string<char, ...
19.21%  a.out    libc.so.6             [.] __mcount_internal
 9.18%  a.out    libstdc++.so.6.0.33   [.] std::__cxx11::basic_string<char, ...
 7.85%  a.out    libc.so.6             [.] isalpha
 4.57%  a.out    libstdc++.so.6.0.33   [.] std::__cxx11::basic_string<char, ...
 4.54%  a.out    libc.so.6             [.] _mcount
 4.33%  a.out    libc.so.6             [.] tolower
 2.75%  a.out    libstdc++.so.6.0.33   [.] std::__cxx11::basic_string<char, ...
 1.99%  a.out    a.out                 [.] bool std::operator==<char, std::char_traits<char>, ...
 1.51%  a.out    a.out                 [.] bool std::operator!=<char, std::char_traits<char>, ...
 1.15%  a.out    libstdc++.so.6.0.33   [.] std::__cxx11::basic_string<char, ...
 1.12%  a.out    libstdc++.so.6.0.33   [.] std::__cxx11::basic_string<char, ...
 0.96%  a.out    libstdc++.so.6.0.33   [.] std::__cxx11::basic_string<char, ...
 0.66%  a.out    a.out                 [.] std::__cxx11::basic_string<char, ...
 0.58%  a.out    libc.so.6             [.] __memmove_avx512_unaligned_erms
\end{minted}

W tym przypadku widzimy, że czas działania jest lepszy o około $30\%$ w przypadku
kompilacji bez optymalizacji. Czas z optymalizacją \verb|-O2| jest identyczny jak
w przypadku poprzedniej wersji co potwierdza tezę, że kompilator mógł zoptymalizować
operacje na iteratorze. \\

Dynamiczne zwiększanie rozmiaru zmiennych \verb|result|, \verb|word| oraz \verb|last_word|
nie jest problematyczne, ponieważ algorytm pracując na dużych plikach tekstowych stosunkowo
szybko osiągnie wymagany rozmiar zmiennych.







\section*{Wersja bare-metal}
W tej wersji użyto tablic znakowych w stylu C o stałym rozmiarze oraz ręcznych operacji
na pamięci z użyciem \verb|new| oraz \verb|memcpy|. Program stał się znacznie mniej 
czytelny oraz wykorzystujemy niebezpieczne operację jak \verb|strcmp|, co wymaga 
znacznie większej ostrożności. Kolejnym problemem jest założenie maksymalnej długości
słowa.

\begin{minted}{cpp}
char* normalize_text(const char* text, const ssize_t size) {
  char* result = new char[size + 1];  // Null-terminated string
  ssize_t result_len = 0;
  char* word = new char[MAX_WORD_LEN];
  ssize_t word_len = 0;
  char* last_word = new char[MAX_WORD_LEN];
  ssize_t last_word_len = 0;
  enum State { NORMAL, COMMA, SPACE, DUPLICATE_SPACE } state = NORMAL;

  for (size_t i = 0; i < size; ++i) {
    char c = text[i];
    if (c < 32 || c > 126) continue;  // rule 1

    if (c == ' ') {
      if (state == SPACE || state == DUPLICATE_SPACE) {  // rule 2
        state = DUPLICATE_SPACE;
      } else {
        state = SPACE;
      }
    } else if (isalpha(c)) {  // rule 3
      c = tolower(c);
      state = NORMAL;
    } else if (ispunct(c)) {  // rule 4
      c = ',';
      state = COMMA;
    } else {
      state = NORMAL;
    }

    if (state == COMMA || state == SPACE) {
      // end of a word
      if (word_len > 0) {  // rule 5
        if (strcmp(word, last_word) != 0) {
          memcpy(result + result_len, word, word_len + 1);
          result_len += word_len;
        }
        memcpy(last_word, word, word_len + 1);
        last_word_len = word_len;
        word_len = 0;
      }
      result[result_len++] = c;  // only delimiter
    } else if (state == NORMAL) {
      // in a word
      word[word_len++] = c;  // only alphanumeric
    }
  }
  if (strcmp(word, last_word) != 0) {  // last word
    memcpy(result + result_len, word, word_len + 1);
    result_len += word_len;
  }

  // Clean up
  result[result_len] = '\0';
  delete[] word;
  delete[] last_word;

  return result;
}
\end{minted}

Czas działania dla plików o rozmiarze \verb|100KB|
\begin{minted}{text}
normalize3 -O0 [Benchmark (1000 iterations): 0.755573 seconds]
normalize3 -O2 [Benchmark (1000 iterations): 0.563578 seconds]
\end{minted}

Wynik profilowania
\begin{minted}{text}
49.18%  a.out    a.out                 [.] normalize_text(char const*, long) 
28.42%  a.out    libc.so.6             [.] __strcmp_evex 
 6.24%  a.out    libc.so.6             [.] tolower 
 1.28%  a.out    libc.so.6             [.] __mcount_internal 
 0.97%  a.out    libc.so.6             [.] _mcount 
 0.44%  a.out    libc.so.6             [.] ispunct 
 0.24%  a.out    a.out                 [.] isalpha@plt 
 0.20%  a.out    a.out                 [.] tolower@plt 
\end{minted}

Tutaj czas działania jest o około $50\%$ krótszy niż w przypadku wersji poprzedniej
pracującej na \verb|string| bez optymalizacji. \\

W przypadku optymalizacji \verb|-O2| czas działania jest już tylko $12\%$ krótszy, 
co także pokazuje, że kompilator dobrze radzi sobię z optymalizacją pracy z napisami
w Cpp. \\

Kod stał się znacznie mniej czytelny oraz bardziej podatny na błędy, co może nie być 
warte takich zmian szczególnie jeśli nie jesteśmny ograniczeni platformą. \\

W wyniku profilowania widzimy, że występują jedyne podstawowe operacje wykorzystywane
w tej wersji algorytmu.






\section*{Wersja równoległa}
Opis wersji. TODO (wymaga rozwiązania problemu łącznia wyników)

\begin{minted}{cpp}
std::string normalize_text_threaded(const std::string& text) {
  unsigned int num_threads = std::thread::hardware_concurrency();

  // For very small texts or single thread systems, use the original function
  if (text.size() < 1000 || num_threads <= 1) {
    return normalize_text(text);
  }

  // Split the text into chunks
  std::vector<std::string> chunks;
  std::vector<std::thread> threads;
  std::vector<std::string> results(num_threads);

  size_t chunk_size = text.size() / num_threads;
  for (unsigned int i = 0; i < num_threads; ++i) {
    size_t start = i * chunk_size;
    size_t end = (i == num_threads - 1) ? text.size() : (i + 1) * chunk_size;
    chunks.push_back(text.substr(start, end - start));
  }

  // Process each chunk in a separate thread
  for (unsigned int i = 0; i < num_threads; ++i) {
    threads.emplace_back(
        [&chunks, &results, i]() { results[i] = normalize_text(chunks[i]); });
  }

  // Wait for all threads to complete
  for (auto& thread : threads) {
    if (thread.joinable()) {
      thread.join();
    }
  }

  std::string combined_result;
  for (const auto& result : results) {
    combined_result += result;
  }
  return combined_result;
}
\end{minted}

Czas działania dla plików o rozmiarze \verb|100KB| przy 24 wątkach
\begin{minted}{text}
normalize4 -O0 [Benchmark (1000 iterations): 1.43878 seconds]
normalize4 -O2 [Benchmark (1000 iterations): 1.25015 seconds] > normalize1 -O2
\end{minted}

Czas działania dla plików o rozmiarze \verb|100MB| przy 24 wątkach
\begin{minted}{text}
normalize4 -O0 [100MB][Benchmark (1 iterations): 0.273202 seconds]
normalize4 -O2 [100MB][Benchmark (1 iterations): 0.196783 seconds]
\end{minted}

Wynik profilowania
\begin{minted}{text}
98.43%  a.out    libc.so.6             [.] __mcount_internal 
 0.79%  a.out    libc.so.6             [.] _mcount 
 0.18%  a.out    a.out                 [.] normalize_text(std::__cxx11::basic_string<char, ...
 0.14%  a.out    libstdc++.so.6.0.33   [.] std::__cxx11::basic_string<char, ...
 0.11%  a.out    libstdc++.so.6.0.33   [.] std::__cxx11::basic_string<char, ...
 0.03%  a.out    libc.so.6             [.] tolower 
 0.03%  a.out    libc.so.6             [.] malloc 
 0.03%  a.out    libc.so.6             [.] __memmove_avx512_unaligned_erms 
 0.02%  a.out    a.out                 [.] __gnu_cxx::__normal_iterator<char const*, ...
 0.02%  a.out    a.out                 [.] __gnu_cxx::__normal_iterator<char const*, ...
 0.02%  a.out    a.out                 [.] bool __gnu_cxx::operator!=<char const*, ...
 0.01%  a.out    libc.so.6             [.] _int_malloc 
 0.01%  a.out    libstdc++.so.6.0.33   [.] std::__cxx11::basic_string<char, ...
 0.01%  a.out    libc.so.6             [.] isalpha 
 0.01%  a.out    libstdc++.so.6.0.33   [.] std::__cxx11::basic_string<char, ...
 0.01%  a.out    libc.so.6             [.] _int_free 
\end{minted}










\section*{Wnioski}
Operacje na ciągach tekstowych są trudne do optymalizacji. 
Nie udało mi się znaleźć miejsca, w którym można by było
wykorzystać wektoryzację. Powodem tego może być dosyć 
skomplikowana logika algorytmu zależna od stanu. \\

Być może rozdzielenie przetwarzania tekstu na kilka przebiegów
umożliwiłoby bardziej efektywne optymalizacje, ale kosztem dodatkowych
iteracji po całości tekstu. \\

Wersja równoległa przyniosła poprawę w przypadku dużych plików (\verb|100MB|) 
gdzie koszt, tworzenia nowych wątków był proprcjonalnie mały do czasu przetwarzania.
Przyśpiesznie wynosiło około $15\%$ na każdy dodatkowy wątek (do 24, 
czyli łącznie $3.6$ razy szybciej), co nie jest idealne. \\

Najważniejszą obserwacją jest to, jak dobrze kompilator radzi sobie
z optymalizacją kodu i implementacją abstrakcji w Cpp.
W porównaniu z nieczytelną wersją w stylu C, wersja podstawowa przy 
użyciu optymalizacji \verb|-O2| była tylko $14\%$ wolniejsza.

\end{document}
